CAGRA fp32/int8, IVF-PQ and IVF-RaBitQ \cite{cagra,pq,ivfrabitq}
use separate parameter grids (\cref{tab:3-grids}); FP32 brute force is
exhaustive. Comparisons match recall within these grids, with memory
reported separately, rather than imposing equal storage or tuning cost. The initial JHQ refinement multipliers are 64, 100, 8 and 4 for
Vogue/arXiv/OpenAI-1536/OpenAI-3072. Additional points extend JHQ's grid
to higher and lower recall; all are retained in the raw records.

\par\medskip\noindent\begin{minipage}{\linewidth}
\centering\captionsetup{hypcap=false}
\captionof{table}[Principal tuning dimensions of the main comparison]{Principal tuning dimensions of the main comparison. Exact per-point parameters and
additional points remain in the raw records. The grid is finite and
differs between methods.}
\label{tab:3-grids}
\footnotesize\setlength{\tabcolsep}{4pt}
\begin{tabularx}{\linewidth}{@{}>{\raggedright\arraybackslash}p{.22\linewidth}>{\raggedright\arraybackslash}X@{}}
\toprule
Family & Recorded search dimensions \\
\midrule
JHQ & Probes and $\alpha$; initial probes 64/128/256 and
$0.5/1/2$ times the workload budget; additional points use up to
1,024 probes and larger/smaller budgets \\
CAGRA fp32/int8 & Graph degree 32/64; search width 1/2/4/8;
intermediate top-$k$ capacity 32--512 \\
IVF-PQ & Eight-bit subcodes; 192/384/768-byte codes at $d=768$,
384/768/1,536 at higher dimensions; fp16/fp32 lookup settings;
probes 16--512 \\
IVF-RaBitQ & 4/8/9 bits per dimension, supported search modes,
probes 16--1,024; additional 4,096-list OpenAI indexes \\
GPU brute force & Exhaustive FP32 squared-L2 search, common host boundary \\
\bottomrule
\end{tabularx}
\end{minipage}\par\medskip

For measured points $a,b$, let $R$ be recall and $Q$ throughput.
The frontier keeps non-dominated points: a point is dominated if another
has no lower recall or throughput and improves at least one:
\[
 a\succ b\iff R_a\geq R_b,\ Q_a\geq Q_b,\
 (R_a,Q_a)\ne(R_b,Q_b).
\]
For target recall $r$, the qualifying measured throughput is
\begin{equation}
Q_{\rm meas}(r)=\max\{Q_i:R_i\geq r\}.
\label{eq:3-measured-frontier}
\end{equation}
Without a qualifying point, this value is undefined. Between adjacent
frontier points bracketing the target, log-throughput interpolation gives
\begin{equation}
\log Q_{\rm int}(r)=(1-\lambda)\log Q_a+\lambda\log Q_b,
\qquad \lambda=\frac{r-R_a}{R_b-R_a}.
\label{eq:3-interpolated-frontier}
\end{equation}
For the interpolated estimate, a measured frontier point within
$5\times10^{-4}$ of the target recall is used directly, taking the first
in increasing-recall order even if it is slightly below the target.
This tolerance applies only to the interpolated estimate; the qualifying
measured throughput above still requires recall at or above the target.
Otherwise, throughput is interpolated on a logarithmic scale between
the adjacent frontier points on either side. Without a nearby point or
a bracket, no interpolated estimate is reported.

Small-batch comparisons also use a qualifying measured endpoint: when
the lowest measured recall already exceeds the target, that point's
throughput is used without extrapolation. Both estimates are reported
when a comparison is close to parity or their rankings differ.
Runs with failed or invalid outputs are excluded from the frontier.
Changes in returned-ID order between repeated calls alone do not
exclude a point (\cref{sec:a-e1-protocol}). CAGRA int8 remains included
(\cref{fig:7-frontiers}).
